\documentclass[10pt]{article}
\usepackage[T1]{fontenc}
\usepackage{lmodern,amsmath,amssymb,amsthm,microtype}
\usepackage[margin=0.82in]{geometry}
\usepackage[colorlinks=true,urlcolor=blue]{hyperref}
\newtheorem{theorem}{Theorem}
\newtheorem{lemma}{Lemma}
\newtheorem{corollary}{Corollary}
\title{A Digit-Sum Congruence for Products of Twin Primes}
\author{Qichao Wang\\{\small Hebei University of Technology}\\[0.3em]
Daoyu Dong\\{\small University of Electronic Science and Technology of China}}
\date{October 26, 2026}
\begin{document}
\maketitle
\begin{abstract}
We prove the congruence conjecture in OEIS A136050: the decimal digit
sum of the product of any twin-prime pair other than $(3,5)$ is
congruent to eight modulo nine. The proof is elementary and applies
to every such pair; it makes no assertion about the infinitude of
twin primes.
\end{abstract}

Let $s(N)$ denote the decimal digit sum of a positive integer $N$.
OEIS A136050 lists $s(p(p+2))$ over the twin-prime pairs in increasing
order, beginning with $s(3\cdot5)=6$. Its comment dated December 2,
2016 conjectures that each subsequent term is one less than a
multiple of nine.

\begin{theorem}
If $p,p+2$ are prime and $p>3$, then
\[
 s(p(p+2))\equiv8\pmod9.
\]
In particular, every term of A136050 except its initial term belongs
to $\{8,17,26,35,\ldots\}$.
\end{theorem}
\begin{proof}
Both primes are odd. Among $p,p+1,p+2$ one number is divisible by
three, and it cannot be either prime since $p>3$. Thus $p+1$ is
divisible by both two and three, giving $p=6m-1$ for an integer
$m\ge1$. Consequently
\[
 p(p+2)=(6m-1)(6m+1)=36m^2-1\equiv-1\pmod9.
\]
For every decimal expansion $N=\sum_j d_j10^j$,
\[
 N\equiv\sum_j d_j=s(N)\pmod9,
\]
which proves the claim. Since the positive digit sum is eight modulo
nine, its least possible value is eight.
\end{proof}

The exceptional pair $(3,5)$ gives $6$, exactly the stated first
term. This settles the entire recorded congruence question. It does
not imply that every value $8+9j$ occurs, nor does it assume or prove
that there are infinitely many twin primes.

The accompanying \texttt{verify.py} compares all 5000 official
b-file terms with a prime-sieve computation of the twin pairs and
with two independent digit-sum implementations. A trial-division
prime enumeration independently checks all 342 twin pairs through
20000. Every one of the 14870 noninitial twin pairs through
2000000 satisfies the theorem. These finite audits check the
implementation; the argument above holds for every pair.

\begin{thebibliography}{1}
\bibitem{oeis} The OEIS Foundation Inc.,
\href{https://oeis.org/A136050}{A136050} and its
\href{https://oeis.org/A136050/b136050.txt}{5000-term b-file},
accessed October 26, 2026.
\end{thebibliography}
\end{document}
